%%%%%%%%%%%%%%%%%%%%%%%%%%%%%%%%%%%%%%%%%%%%%%%%%%%%%%%%%%%%%%%%%%%%%%
% Lớp 4 | Chủ đề 1. Ôn tập và bổ sung | Giải bài toán có ba bước tính
% Dạng 4D1?5-3: Giả thiết tạm (gà – chó, tiền xu)
%%%%%%%%%%%%%%%%%%%%%%%%%%%%%%%%%%%%%%%%%%%%%%%%%%%%%%%%%%%%%%%%%%%%%%


% nguồn: 2026-11-18_lop4_chu-de-1-tu-luan (nhập 18/11/2026)
\begin{bt}%[4D1H5-3]%[ma:34a5bad9]
% nguon: Toán tư duy lớp 4 – Bài 5 SGK Toán 4 KNTT
	Vừa gà vừa chó có $36$ con, đếm được $100$ cái chân. Hỏi có bao nhiêu con gà, bao nhiêu con chó?
	\loigiai{
		Giả sử cả $36$ con đều là gà thì có $36 \times 2 = 72$ (chân).\\
		Số chân bị thiếu: $100 - 72 = 28$ (chân).\\
		Mỗi con chó có nhiều hơn con gà $2$ chân, nên số chó: $28 : 2 = 14$ (con).\\
		Số gà: $36 - 14 = 22$ (con).\\
		Thử lại: $22 \times 2 + 14 \times 4 = 44 + 56 = 100$ (chân).
	}
\end{bt}

% nguồn: 2026-11-18_lop4_chu-de-1-tu-luan (nhập 18/11/2026)
\begin{bt}%[4D1V5-3]%[ma:26fa2b9f]
% nguon: Toán tư duy lớp 4 – Bài 5 SGK Toán 4 KNTT
	Minh có $15$ tờ tiền gồm hai loại $2\,000$ đồng và $5\,000$ đồng, tổng cộng là $54\,000$ đồng. Hỏi Minh có mấy tờ mỗi loại?
	\loigiai{
		Giả sử cả $15$ tờ đều là loại $2\,000$ đồng thì có $2\,000 \times 15 = 30\,000$ (đồng).\\
		Số tiền còn thiếu: $54\,000 - 30\,000 = 24\,000$ (đồng).\\
		Mỗi tờ $5\,000$ đồng hơn tờ $2\,000$ đồng là $3\,000$ đồng, nên số tờ $5\,000$ đồng: $24\,000 : 3\,000 = 8$ (tờ).\\
		Số tờ $2\,000$ đồng: $15 - 8 = 7$ (tờ).\\
		Thử lại: $7 \times 2\,000 + 8 \times 5\,000 = 14\,000 + 40\,000 = 54\,000$ (đồng).
	}
\end{bt}

% nguồn: 2026-11-18_lop4_chu-de-1-tu-luan (nhập 18/11/2026)
\begin{bt}%[4D1C5-3]%[ma:08eafc21]
% nguon: Toán tư duy lớp 4 – Bài 5 SGK Toán 4 KNTT
	Một đề thi có $20$ câu. Mỗi câu làm đúng được $5$ điểm, mỗi câu làm sai hoặc bỏ trống bị trừ $2$ điểm. Bạn Hà được $65$ điểm. Hỏi Hà làm đúng bao nhiêu câu?
	\loigiai{
		Giả sử Hà làm đúng cả $20$ câu thì được $5 \times 20 = 100$ (điểm).\\
		Hà bị hụt: $100 - 65 = 35$ (điểm).\\
		Mỗi câu sai thay vì đúng thì mất $5$ điểm và bị trừ thêm $2$ điểm, tức là hụt $5 + 2 = 7$ (điểm).\\
		Số câu sai: $35 : 7 = 5$ (câu). Số câu đúng: $20 - 5 = 15$ (câu).\\
		Thử lại: $15 \times 5 - 5 \times 2 = 75 - 10 = 65$ (điểm).
	}
\end{bt}

% nguồn: 2026-11-18_lop4_chu-de-1-tu-luan-bo-sung (nhập 18/11/2026)
\begin{bt}%[4D1V5-3]%[ma:f572cb93]
% nguon: Toán tư duy lớp 4 – Bài 5 SGK Toán 4 KNTT
	Lớp em có $32$ bạn đi tham quan bằng $8$ chiếc thuyền gồm thuyền lớn và thuyền nhỏ. Mỗi thuyền lớn chở $5$ bạn, mỗi thuyền nhỏ chở $3$ bạn, các thuyền đều chở vừa đủ. Hỏi có mấy thuyền mỗi loại?
	\loigiai{
		Giả sử cả $8$ thuyền đều là thuyền nhỏ thì chở được $3 \times 8 = 24$ (bạn).\\
		Còn thiếu chỗ cho: $32 - 24 = 8$ (bạn).\\
		Mỗi thuyền lớn chở hơn thuyền nhỏ $5 - 3 = 2$ (bạn), nên số thuyền lớn: $8 : 2 = 4$ (thuyền).\\
		Số thuyền nhỏ: $8 - 4 = 4$ (thuyền).\\
		Thử lại: $4 \times 5 + 4 \times 3 = 20 + 12 = 32$ (bạn).
	}
\end{bt}
